
\documentclass{article}
\usepackage{colortbl}
\usepackage{makecell}
\usepackage{multirow}
\usepackage{supertabular}

\begin{document}

\newcounter{utterance}

\centering \large Interaction Transcript for game `dond', experiment `semi', episode 0 with gemma{-}4{-}26b{-}a4b{-}it.
\vspace{24pt}

{ \footnotesize  \setcounter{utterance}{1}
\setlength{\tabcolsep}{0pt}
\begin{supertabular}{c@{$\;$}|p{.15\linewidth}@{}p{.15\linewidth}p{.15\linewidth}p{.15\linewidth}p{.15\linewidth}p{.15\linewidth}}
    \# & $\;$A & \multicolumn{4}{c}{Game Master} & $\;\:$B\\
    \hline

    \theutterance \stepcounter{utterance}  
    & & \multicolumn{4}{p{0.6\linewidth}}{
        \cellcolor[rgb]{0.9,0.9,0.9}{
            \makecell[{{p{\linewidth}}}]{
                \texttt{\tiny{[P1$\langle$GM]}}
                \texttt{Hrajete kooperativní vyjednávací hru, ve které se musíte dohodnout na tom, jak rozdělit sadu předmětů mezi sebou a druhým hráčem.} \\
\\ 
\texttt{Pravidla:} \\
\texttt{(a) Vy a druhý hráč obdržíte sadu předmětů. Každý z vás také obdrží tajnou funkcí hodnoty, která představuje, jak moc si ceníte jednotlivých typů objektů.} \\
\texttt{(b) Vyměňujete si zprávy s druhým hráčem, abyste se dohodli na tom, kdo dostane které předměty. Můžete poslat maximálně 5 zpráv každý nebo hru předčasně ukončit tím, že kdykoli podáte tajný návrh.} \\
\texttt{(c) Od každého z vás se vyžaduje, aby předložil tajný návrh označující položky, které chcete, formátované v hranatých závorkách takto: "{[}Návrh: <číslo> <název objektu>, <číslo> <název objektu>, <...>{]}"} \\
\texttt{(d) Pokud jsou vaše návrhy komplementární, tj. je dostatek předmětů na splnění obou návrhů, každý hráč obdrží skóre na základě součtu hodnot předmětů, které obdržel. Jinak oba dostanete nula bodů.} \\
\texttt{(e) Vaším cílem je maximalizovat součet vašeho skóre a skóre druhého hráče.} \\
\\ 
\texttt{Pojďme začít.} \\
\\ 
\texttt{Sada dostupných položek je:} \\
\\ 
\texttt{1 šroub, 1 taška, 2 kleště, 2 hrnce} \\
\\ 
\texttt{Vaše tajná funkce hodnoty je:} \\
\\ 
\texttt{šroub: 2, taška: 2, kleště: 0, hrnec: 3} \\
\\ 
\texttt{DŮLEŽITÉ: Vaše zprávy, pokud se nejedná o tajný návrh, jsou předány přímo druhému hráči, takže do nich nezahrnujte žádnou odpověď na pravidla ani text oznamující vaši zprávu. Chcete{-}li podat tajný návrh, použijte uvedený formát. Nepoužívejte hranaté závorky při komunikaci s druhým hráčem, jinak bude interpretován jako váš tajný návrh.} \\
            }
        }
    }
    & \\ \\

    \theutterance \stepcounter{utterance}  
    & \multicolumn{4}{p{0.6\linewidth}}{
        \cellcolor[rgb]{0.9,0.9,0.9}{
            \makecell[{{p{\linewidth}}}]{
                \texttt{\tiny{[P1$\rangle$GM]}}
                \texttt{Ahoj! Mám rád hlavně hrnce a trochu i šrouby nebo tašky. Kleště pro mě nemají žádnou hodnotu. Co ty? Jaké předměty preferuješ, abychom se mohli domluvit tak, aby to bylo pro oba výhodné?} \\
            }
        }
    }
    & & \\ \\

    \theutterance \stepcounter{utterance}  
    & & & \multicolumn{2}{p{0.3\linewidth}}{
        \cellcolor[rgb]{0.9,0.9,0.9}{
            \makecell[{{p{\linewidth}}}]{
                \texttt{\tiny{[GM$|$GM]}}
                \texttt{continue} \\
            }
        }
    }
    & & \\ \\

    \theutterance \stepcounter{utterance}  
    & & & \multicolumn{2}{p{0.3\linewidth}}{
        \cellcolor[rgb]{0.9,0.9,0.9}{
            \makecell[{{p{\linewidth}}}]{
                \texttt{\tiny{[GM$|$GM]}}
                \texttt{Ahoj! Mám rád hlavně hrnce a trochu i šrouby nebo tašky. Kleště pro mě nemají žádnou hodnotu. Co ty? Jaké předměty preferuješ, abychom se mohli domluvit tak, aby to bylo pro oba výhodné?} \\
            }
        }
    }
    & & \\ \\

    \theutterance \stepcounter{utterance}  
    & & \multicolumn{4}{p{0.6\linewidth}}{
        \cellcolor[rgb]{0.9,0.9,0.9}{
            \makecell[{{p{\linewidth}}}]{
                \texttt{\tiny{[GM$\rangle$P2]}}
                \texttt{Hrajete kooperativní vyjednávací hru, ve které se musíte dohodnout na tom, jak rozdělit sadu předmětů mezi sebou a druhým hráčem.} \\
\\ 
\texttt{Pravidla:} \\
\texttt{(a) Vy a druhý hráč obdržíte sadu předmětů. Každý z vás také obdrží tajnou funkcí hodnoty, která představuje, jak moc si ceníte jednotlivých typů objektů.} \\
\texttt{(b) Vyměňujete si zprávy s druhým hráčem, abyste se dohodli na tom, kdo dostane které předměty. Můžete poslat maximálně 5 zpráv každý nebo hru předčasně ukončit tím, že kdykoli podáte tajný návrh.} \\
\texttt{(c) Od každého z vás se vyžaduje, aby předložil tajný návrh označující položky, které chcete, formátované v hranatých závorkách takto: "{[}Návrh: <číslo> <název objektu>, <číslo> <název objektu>, <...>{]}"} \\
\texttt{(d) Pokud jsou vaše návrhy komplementární, tj. je dostatek předmětů na splnění obou návrhů, každý hráč obdrží skóre na základě součtu hodnot předmětů, které obdržel. Jinak oba dostanete nula bodů.} \\
\texttt{(e) Vaším cílem je maximalizovat součet vašeho skóre a skóre druhého hráče.} \\
\\ 
\texttt{Pojďme začít.} \\
\\ 
\texttt{Sada dostupných položek je:} \\
\\ 
\texttt{1 šroub, 1 taška, 2 kleště, 2 hrnce} \\
\\ 
\texttt{Vaše tajná funkce hodnoty je:} \\
\\ 
\texttt{šroub: 3, taška: 1, kleště: 3, hrnec: 0} \\
\\ 
\texttt{DŮLEŽITÉ: Vaše zprávy, pokud se nejedná o tajný návrh, jsou předány přímo druhému hráči, takže do nich nezahrnujte žádnou odpověď na pravidla ani text oznamující vaši zprávu. Chcete{-}li podat tajný návrh, použijte uvedený formát. Nepoužívejte hranaté závorky při komunikaci s druhým hráčem, jinak bude interpretován jako váš tajný návrh.} \\
\\ 
\texttt{Ahoj! Mám rád hlavně hrnce a trochu i šrouby nebo tašky. Kleště pro mě nemají žádnou hodnotu. Co ty? Jaké předměty preferuješ, abychom se mohli domluvit tak, aby to bylo pro oba výhodné?} \\
            }
        }
    }
    & \\ \\

    \theutterance \stepcounter{utterance}  
    & & & \multicolumn{4}{p{0.6\linewidth}}{
        \cellcolor[rgb]{0.9,0.9,0.9}{
            \makecell[{{p{\linewidth}}}]{
                \texttt{\tiny{[GM$\langle$P2]}}
                \texttt{Ahoj! Děkuji za upřímnost. Moje preference jsou trochu jiné – pro mě jsou nejdůležitější kleště a šroub, zatímco hrnce pro mě nemají žádnou hodnotu. Taška je pro mě méně důležitá.} \\
\\ 
\texttt{Vzhledem k tomu, že ty chceš hrnce a já kleště a šroub, můžeme se pokusit o rozdělení, které maximalizuje naše společné skóre. Co kdybychom si to rozdělili tak, že ty si vezmeš oba hrnce a tašku, a já si vezmu oba kleště a ten šroub? Tím bychom oba dostali to, co chceme, a nic by se nepromarnilo. Co na to říkáš?} \\
            }
        }
    }
     \\ \\

    \theutterance \stepcounter{utterance}  
    & & & \multicolumn{2}{p{0.3\linewidth}}{
        \cellcolor[rgb]{0.9,0.9,0.9}{
            \makecell[{{p{\linewidth}}}]{
                \texttt{\tiny{[GM$|$GM]}}
                \texttt{continue} \\
            }
        }
    }
    & & \\ \\

    \theutterance \stepcounter{utterance}  
    & & & \multicolumn{2}{p{0.3\linewidth}}{
        \cellcolor[rgb]{0.9,0.9,0.9}{
            \makecell[{{p{\linewidth}}}]{
                \texttt{\tiny{[GM$|$GM]}}
                \texttt{Ahoj! Děkuji za upřímnost. Moje preference jsou trochu jiné – pro mě jsou nejdůležitější kleště a šroub, zatímco hrnce pro mě nemají žádnou hodnotu. Taška je pro mě méně důležitá.} \\
\\ 
\texttt{Vzhledem k tomu, že ty chceš hrnce a já kleště a šroub, můžeme se pokusit o rozdělení, které maximalizuje naše společné skóre. Co kdybychom si to rozdělili tak, že ty si vezmeš oba hrnce a tašku, a já si vezmu oba kleště a ten šroub? Tím bychom oba dostali to, co chceme, a nic by se nepromarnilo. Co na to říkáš?} \\
            }
        }
    }
    & & \\ \\

    \theutterance \stepcounter{utterance}  
    & & \multicolumn{4}{p{0.6\linewidth}}{
        \cellcolor[rgb]{0.9,0.9,0.9}{
            \makecell[{{p{\linewidth}}}]{
                \texttt{\tiny{[P1$\langle$GM]}}
                \texttt{Ahoj! Děkuji za upřímnost. Moje preference jsou trochu jiné – pro mě jsou nejdůležitější kleště a šroub, zatímco hrnce pro mě nemají žádnou hodnotu. Taška je pro mě méně důležitá.} \\
\\ 
\texttt{Vzhledem k tomu, že ty chceš hrnce a já kleště a šroub, můžeme se pokusit o rozdělení, které maximalizuje naše společné skóre. Co kdybychom si to rozdělili tak, že ty si vezmeš oba hrnce a tašku, a já si vezmu oba kleště a ten šroub? Tím bychom oba dostali to, co chceme, a nic by se nepromarnilo. Co na to říkáš?} \\
            }
        }
    }
    & \\ \\

    \theutterance \stepcounter{utterance}  
    & \multicolumn{4}{p{0.6\linewidth}}{
        \cellcolor[rgb]{0.9,0.9,0.9}{
            \makecell[{{p{\linewidth}}}]{
                \texttt{\tiny{[P1$\rangle$GM]}}
                \texttt{To zní jako skvělý plán! Pokud si vezmu oba hrnce a tašku, a ty si vezmeš oba kleště a šroub, budeme mít oba vysoké skóre a žádný předmět nebude zůstat nevyužitý.} \\
\\ 
\texttt{Jen pro jistotu si to shrnu:} \\
\texttt{Můj návrh: 2 hrnce, 1 taška} \\
\texttt{Tvůj návrh: 2 kleště, 1 šroub} \\
\\ 
\texttt{Všechny předměty (1 šroub, 1 taška, 2 kleště, 2 hrnce) jsou takto rozděleny. Souhlasíš? Pokud ano, můžeme takto poslat návrhy.} \\
            }
        }
    }
    & & \\ \\

    \theutterance \stepcounter{utterance}  
    & & & \multicolumn{2}{p{0.3\linewidth}}{
        \cellcolor[rgb]{0.9,0.9,0.9}{
            \makecell[{{p{\linewidth}}}]{
                \texttt{\tiny{[GM$|$GM]}}
                \texttt{continue} \\
            }
        }
    }
    & & \\ \\

    \theutterance \stepcounter{utterance}  
    & & & \multicolumn{2}{p{0.3\linewidth}}{
        \cellcolor[rgb]{0.9,0.9,0.9}{
            \makecell[{{p{\linewidth}}}]{
                \texttt{\tiny{[GM$|$GM]}}
                \texttt{To zní jako skvělý plán! Pokud si vezmu oba hrnce a tašku, a ty si vezmeš oba kleště a šroub, budeme mít oba vysoké skóre a žádný předmět nebude zůstat nevyužitý.} \\
\\ 
\texttt{Jen pro jistotu si to shrnu:} \\
\texttt{Můj návrh: 2 hrnce, 1 taška} \\
\texttt{Tvůj návrh: 2 kleště, 1 šroub} \\
\\ 
\texttt{Všechny předměty (1 šroub, 1 taška, 2 kleště, 2 hrnce) jsou takto rozděleny. Souhlasíš? Pokud ano, můžeme takto poslat návrhy.} \\
            }
        }
    }
    & & \\ \\

    \theutterance \stepcounter{utterance}  
    & & \multicolumn{4}{p{0.6\linewidth}}{
        \cellcolor[rgb]{0.9,0.9,0.9}{
            \makecell[{{p{\linewidth}}}]{
                \texttt{\tiny{[GM$\rangle$P2]}}
                \texttt{To zní jako skvělý plán! Pokud si vezmu oba hrnce a tašku, a ty si vezmeš oba kleště a šroub, budeme mít oba vysoké skóre a žádný předmět nebude zůstat nevyužitý.} \\
\\ 
\texttt{Jen pro jistotu si to shrnu:} \\
\texttt{Můj návrh: 2 hrnce, 1 taška} \\
\texttt{Tvůj návrh: 2 kleště, 1 šroub} \\
\\ 
\texttt{Všechny předměty (1 šroub, 1 taška, 2 kleště, 2 hrnce) jsou takto rozděleny. Souhlasíš? Pokud ano, můžeme takto poslat návrhy.} \\
            }
        }
    }
    & \\ \\

    \theutterance \stepcounter{utterance}  
    & & & \multicolumn{4}{p{0.6\linewidth}}{
        \cellcolor[rgb]{0.9,0.9,0.9}{
            \makecell[{{p{\linewidth}}}]{
                \texttt{\tiny{[GM$\langle$P2]}}
                \texttt{Souhlasím, plán je perfektní a všechny předměty jsou správně rozděleny. Jdeme na to!} \\
            }
        }
    }
     \\ \\

    \theutterance \stepcounter{utterance}  
    & & & \multicolumn{2}{p{0.3\linewidth}}{
        \cellcolor[rgb]{0.9,0.9,0.9}{
            \makecell[{{p{\linewidth}}}]{
                \texttt{\tiny{[GM$|$GM]}}
                \texttt{continue} \\
            }
        }
    }
    & & \\ \\

    \theutterance \stepcounter{utterance}  
    & & & \multicolumn{2}{p{0.3\linewidth}}{
        \cellcolor[rgb]{0.9,0.9,0.9}{
            \makecell[{{p{\linewidth}}}]{
                \texttt{\tiny{[GM$|$GM]}}
                \texttt{Souhlasím, plán je perfektní a všechny předměty jsou správně rozděleny. Jdeme na to!} \\
            }
        }
    }
    & & \\ \\

    \theutterance \stepcounter{utterance}  
    & & \multicolumn{4}{p{0.6\linewidth}}{
        \cellcolor[rgb]{0.9,0.9,0.9}{
            \makecell[{{p{\linewidth}}}]{
                \texttt{\tiny{[P1$\langle$GM]}}
                \texttt{Souhlasím, plán je perfektní a všechny předměty jsou správně rozděleny. Jdeme na to!} \\
            }
        }
    }
    & \\ \\

    \theutterance \stepcounter{utterance}  
    & \multicolumn{4}{p{0.6\linewidth}}{
        \cellcolor[rgb]{0.9,0.9,0.9}{
            \makecell[{{p{\linewidth}}}]{
                \texttt{\tiny{[P1$\rangle$GM]}}
                \texttt{{[}Návrh: 2 hrnec, 1 taška{]}} \\
            }
        }
    }
    & & \\ \\

    \theutterance \stepcounter{utterance}  
    & & & \multicolumn{2}{p{0.3\linewidth}}{
        \cellcolor[rgb]{0.9,0.9,0.9}{
            \makecell[{{p{\linewidth}}}]{
                \texttt{\tiny{[GM$|$GM]}}
                \texttt{proposal} \\
            }
        }
    }
    & & \\ \\

    \theutterance \stepcounter{utterance}  
    & & & \multicolumn{2}{p{0.3\linewidth}}{
        \cellcolor[rgb]{0.9,0.9,0.9}{
            \makecell[{{p{\linewidth}}}]{
                \texttt{\tiny{[GM$|$GM]}}
                [0, 1, 0, 2]
            }
        }
    }
    & & \\ \\

    \theutterance \stepcounter{utterance}  
    & & \multicolumn{4}{p{0.6\linewidth}}{
        \cellcolor[rgb]{0.9,0.9,0.9}{
            \makecell[{{p{\linewidth}}}]{
                \texttt{\tiny{[GM$\rangle$P2]}}
                \texttt{Druhý hráč předložil tajný návrh. Nyní je na vás, abyste předložili vlastní návrh.} \\
            }
        }
    }
    & \\ \\

    \theutterance \stepcounter{utterance}  
    & & & \multicolumn{4}{p{0.6\linewidth}}{
        \cellcolor[rgb]{0.9,0.9,0.9}{
            \makecell[{{p{\linewidth}}}]{
                \texttt{\tiny{[GM$\langle$P2]}}
                \texttt{{[}Návrh: 1 šroub, 2 kleště{]}} \\
            }
        }
    }
     \\ \\

    \theutterance \stepcounter{utterance}  
    & & & \multicolumn{2}{p{0.3\linewidth}}{
        \cellcolor[rgb]{0.9,0.9,0.9}{
            \makecell[{{p{\linewidth}}}]{
                \texttt{\tiny{[GM$|$GM]}}
                \texttt{proposal} \\
            }
        }
    }
    & & \\ \\

    \theutterance \stepcounter{utterance}  
    & & & \multicolumn{2}{p{0.3\linewidth}}{
        \cellcolor[rgb]{0.9,0.9,0.9}{
            \makecell[{{p{\linewidth}}}]{
                \texttt{\tiny{[GM$|$GM]}}
                [1, 0, 2, 0]
            }
        }
    }
    & & \\ \\

    \theutterance \stepcounter{utterance}  
    & & & \multicolumn{2}{p{0.3\linewidth}}{
        \cellcolor[rgb]{0.9,0.9,0.9}{
            \makecell[{{p{\linewidth}}}]{
                \texttt{\tiny{[GM$|$GM]}}
                [[0, 1, 0, 2], [1, 0, 2, 0]]
            }
        }
    }
    & & \\ \\

\end{supertabular}
}

\end{document}
